The moving-knife algorithm (MKA) is a structure algorithm in \textsc{bisect}
(\texttt{--structure moving-knife}) designed to choose directional recursive
splits with a strong worst-side compactness proxy.

\subsection{Algorithm Description}

MKA adapts the moving-knife intuition~\citep{brams2000win} to a recursive
tract graph splitter.  At each recursive step, the implementation tests a grid
of candidate directions, projects tract centroids onto each direction, sorts
units by projection rank, and cuts the first population-balanced prefix.  It
then scores the two sides with a point-cloud enclosing-circle proxy and selects
the direction with the best worst-side score before bounded local rebalance.

Formally, at step $t$ of the recursion with region $G_t$ to be divided into
$k_1$ and $k_2$ districts:
\begin{enumerate}
  \item Generate candidate directions $\{\theta_1, \ldots, \theta_m\}$.
  \item For each direction, project centroids, sort units, and form a
    population-balanced prefix/suffix cut.
  \item Compute a proxy score
    $\text{ReockProxy}^{\min}(\theta_j) =
    \min(\widehat R(L_j), \widehat R(R_j))$ over tract-centroid point clouds.
  \item Select $\theta^* = \arg\max_j \text{ReockProxy}^{\min}(\theta_j)$,
    then run the local rebalance path.
\end{enumerate}

\subsection{Reock Maximisation Property}

MKA does not guarantee a globally maximum Reock partition and does not optimise
the exact polygon-MBC metric.  It provides an auditable directional baseline
whose candidate ledger can be checked against the implemented point-cloud proxy.

\subsection{Empirical Validation}

Table~\ref{tab:mka_reock} reports mean and minimum district Reock-proxy values
for all four algorithms on NC 2020 (seed $= 42^\dagger$).

\begin{table}[ht]
\centering
\caption{District reported Reock under four structure algorithms, NC 2020, seed $= 42^\dagger$.}
\label{tab:mka_reock}
\begin{tabular}{lcccc}
\toprule
Algorithm & Min & Median & Mean & Max \\
\midrule
standard-bisect & 0.21 & 0.29 & 0.31$^\dagger$ & 0.46 \\
prime-factor    & 0.24 & 0.32 & 0.34$^\dagger$ & 0.49 \\
ratio-optimal   & 0.26 & 0.35 & 0.37$^\dagger$ & 0.51 \\
moving-knife    & 0.29 & 0.37 & 0.38$^\dagger$ & 0.53 \\
\bottomrule
\end{tabular}
\begin{tablenotes}
\footnotesize
\item $^\dagger$Single-run result, seed $= 42$.
\end{tablenotes}
\end{table}

MKA achieves mean reported Reock $= 0.38^\dagger$ on NC 2020---a 23\% improvement over
standard-bisect ($0.31^\dagger$) and a 3\% improvement over ratio-optimal ($0.37^\dagger$).
The minimum district Reock under MKA is $0.29$, versus $0.21$ for standard-bisect
(a 38\% improvement in the worst-case district).
This supports the weaker implementation claim: MKA's directional sweep is aligned
with its worst-side compactness proxy on this single-run fixture.  Stronger
same-metric or statewide superiority claims require archived per-node candidate
ledgers and exact polygon-Reock comparisons.
